\documentclass[conference]{IEEEtran}

% T1 fontenc: gives Times (ptm) real small-caps coverage for the glyphs that
% otherwise fall back to TS1 inside IEEEtran's small-caps table captions
% (source of the "Font shape TS1/ptm/m/sc undefined" warning).
\usepackage[T1]{fontenc}
\usepackage{cite}
\usepackage{amsmath,amssymb,amsfonts}
\usepackage{graphicx}
\usepackage{textcomp}
\usepackage{xcolor}
\usepackage{booktabs}
\usepackage{url}
\usepackage{array}
\usepackage{tikz}
\usetikzlibrary{arrows.meta,positioning,fit,backgrounds,calc,shapes.geometric,shapes.misc}

\begin{document}

\title{Cluster-Routed Hard-Negative Mining for UAV Cross-View Visual-Localization}

\author{
\IEEEauthorblockN{Kim-Phuong Phung}
\IEEEauthorblockA{
% TODO: department/faculty name and email
\textit{Le Quy Don Technical University}\\
Hanoi, Vietnam\\
email@example.com}
\and
\IEEEauthorblockN{Quang-Uy Nguyen}
\IEEEauthorblockA{
% TODO: department/faculty name and email
\textit{Le Quy Don Technical University}\\
Hanoi, Vietnam\\
email@example.com}
}

\maketitle

\begin{abstract}
UAV visual localization retrieves the geographic position of an aerial
query image from a large georeferenced satellite gallery. Flat contrastive
training draws negatives uniformly from the batch or dataset, so the
gradient concentrates on already-easy pairs while the visually repetitive
candidates that actually cause retrieval errors -- nearby rooftops, road
grids, field boundaries -- are rarely contrasted against the query. We
propose a cluster-routed pipeline in which a single raw, L2-normalized
backbone embedding serves as the retrieval descriptor with no projection
head, trained so that a K-means partition over this descriptor doubles as
both the evaluation-time routing structure and the training-time
hard-negative source: negatives are drawn from a query's own cluster and
its nearest neighbors, ranked by descriptor similarity, with co-positive
exclusion preventing true matches from being pushed apart. Because
clusters are rebuilt from the current descriptor, the mining pool tracks
the model and continually re-targets whatever confusions remain. On five
UAV-satellite benchmarks (DenseUAV, SUES-200, University-1652, GTA-UAV,
UAV-VisLoc) this single mechanism matches or exceeds recent
part-supervised and domain-alignment methods that require additional
heads or losses.
\end{abstract}

\begin{IEEEkeywords}
UAV localization, cross-view retrieval, hard-negative mining, clustering,
metric learning.
\end{IEEEkeywords}

\section{Introduction}
UAV visual localization matches an aerial query against a georeferenced
satellite gallery, learning an embedding in which geographically
corresponding views are close and unrelated locations are far apart
\cite{hadsell2006dimensionality,schroff2015facenet}. The central obstacle
is not the appearance gap between views but the many visually repetitive
candidates a large map contains: a flat in-batch contrastive objective
draws negatives uniformly from the batch, so most land far from the
query, where the gradient is already near zero, while the query's actual
look-alike neighbors are almost never sampled.

We address this with a single mechanism rather than an additional loss
term or head. The backbone's own L2-normalized output is trained -- by
unfreezing its last stages -- to place corresponding query-gallery crops
into the same K-means cluster. This descriptor is then used for three
roles at once: routing a query to its nearest map clusters at evaluation
time, scoring candidates inside those clusters, and, centrally,
\emph{supplying the hard-negative pool during training}: negatives are
drawn from a query's own cluster and its nearest neighbors, ranked by
current descriptor similarity, with co-positive exclusion removing any
candidate that shares a positive with the query. Because the partition is
rebuilt from the evolving descriptor, the negative pool automatically
re-targets whatever confusions the model has not yet resolved.

Our contribution is this cluster-structured hard-negative mining
mechanism and the coarse-to-fine retrieval it enables. We evaluate it on
five UAV-satellite benchmarks against recent methods that instead add
part-partitioning heads or explicit spatial-alignment losses, and show
the single-descriptor mechanism matches or exceeds them without extra
trainable parameters.

\section{Related Work}
Sample4Geo showed that symmetric InfoNCE with hard-negative sampling
outperforms complex aggregation pipelines for cross-view retrieval
\cite{deuser2023sample4geo}, and Game4Loc extended this to a
partial-overlap UAV setting with weighted InfoNCE over a multi-zoom
satellite tile gallery \cite{ji2024game4loc}. A parallel line strengthens
this baseline with explicit spatial supervision instead of sampling
changes: LPN and FSRA partition the feature map into rings or
activation-ranked regions and supervise each part
\cite{wang2021lpn,dai2022fsra}; MCCG applies part-level contrastive
matching \cite{shen2024mccg}; DAC adds a domain/scene spatial-alignment
loss and hard-negative reweighting \cite{xia2024dac}; MEAN extends this
with a multilevel consistency loss \cite{chen2025mean}; CAMP combines
position-aware partitioning with attribute mining \cite{wu2024camp}; and
CEUSP augments a compact backbone with multi-directional attention and
dynamic sampling \cite{xu2025ceusp}. GLVL instead separates coarse
retrieval from a dedicated fine-matching module \cite{li2023glvl}. All of
these add parameters or losses beyond the pooled descriptor; the method
here instead restructures \emph{what the existing descriptor is trained
against}, using FAISS K-means \cite{johnson2019billion,macqueen1967methods}
as both the routing and negative-sampling structure.

\section{Method}
\subsection{Descriptor and Cluster Routing}
The visual encoder is a pretrained \texttt{timm} backbone (Swin-B at
$384{\times}384$ \cite{liu2021swin}, or a $3\times$-smaller ConvNeXt-T/B
variant) with the classification layer removed. Its globally pooled
output is L2-normalized directly into the retrieval descriptor,
\begin{equation}
    \mathbf{c}(\mathbf{x}) = \mathrm{norm}\big(f_b(\mathbf{x})\big)
    \in \mathbb{R}^{D},
\end{equation}
with no projection head: the same $\mathbf{c}(\mathbf{x})$ is used for
matching, K-means assignment, and negative mining. Only the last backbone
stages are unfrozen, so this is the sole mechanism by which the
descriptor is adapted. Given L2-normalized centroids
$\mathbf{C}=\{\mathbf{c}_1,\dots,\mathbf{c}_K\}$, the routing score is
$\mathbf{s}(\mathbf{x})=\mathbf{c}(\mathbf{x})\mathbf{C}^T$; at evaluation
time a query is routed to its top-$K_c$ nearest clusters and only the
tiles assigned to them are scored, cutting candidate-scoring cost by a
factor of $K/K_c$.

\subsection{Cluster-Structured Hard-Negative Mining}
\label{sec:hard_mining}
At the start of each epoch the trainer extracts descriptors for all
anchors and positives, runs FAISS K-means, and measures the
anchor-positive same-cluster probability, re-clustering (warm-started
from the previous centroids after the first epoch) whenever it drops
below 0.80. This partition then supplies training-time negatives: for
each anchor, 50\% of negatives are drawn from its own cluster (ranked by
descriptor similarity to the anchor -- the hardest confusers), 30\% from
the nearest neighboring cluster, and 20\% globally at random for
coverage. Two safeguards keep this signal clean. \emph{Co-positive
exclusion} removes any candidate sharing a positive tile with the anchor
(directly or through an overlapping positive) from the negative pool,
since pushing apart two views of the same place would oppose the
alignment objective. The negative loss is also warmed up over early
training steps so unstable early descriptors are not dominated by noisy
hard negatives. Because clusters are rebuilt from the current descriptor,
this pool is self-correcting: as look-alike regions are pulled apart, the
clusters -- and therefore the negatives -- re-focus on whatever confusion
remains, unlike a fixed or randomly-refreshed negative set.
Fig.~\ref{fig:mining_comparison} contrasts this against standard in-batch
InfoNCE sampling.

\begin{figure}[t]
\centering
\tikzset{
  negopen/.style = {circle, draw=red!70!black, thick, inner sep=1.4pt},
  negfill/.style = {circle, fill=red!70!black, inner sep=1.4pt},
  posdot/.style  = {circle, fill=green!55!black, inner sep=1.6pt},
  anchdot/.style = {star, star points=5, star point ratio=2.4, fill=black,
                    inner sep=1.3pt},
  graydot/.style = {circle, fill=black!30, inner sep=1.2pt},
  cent/.style    = {cross out, draw=blue!60!black, thick, inner sep=1.6pt},
  panel/.style   = {rounded corners=4pt, draw=black!60, thick},
  note/.style    = {font=\tiny, align=center},
}
\resizebox{\columnwidth}{!}{%
\begin{tikzpicture}

\begin{scope}
\draw[panel] (0,0) rectangle (7,5);
\node[font=\scriptsize\bfseries, anchor=south west] at (0.05,5.05)
  {(a) In-batch InfoNCE};
\node[anchdot] (a1) at (2.2,2.6) {};
\node[posdot]  at (2.55,2.95) {};
\draw[dashed, black!50] (2.2,2.6) circle (0.85);
\node[note, black!60] at (2.2,1.35) {look-alike neighborhood:\\rarely sampled};
\foreach \p in {(5.9,4.3),(6.2,1.2),(4.8,0.7),(5.4,3.1),(6.5,2.5),
                (4.3,4.5),(3.9,1.0),(1.0,4.4),(0.8,0.9)}
  \node[negopen] at \p {};
\node[note] at (4.9,4.85) {negatives $=$ other batch items};
\node[note, red!60!black] at (5.6,0.35) {distant $\Rightarrow$ easy};
\end{scope}

\begin{scope}[shift={(7.6,0)}]
\draw[panel] (0,0) rectangle (7,5);
\node[font=\scriptsize\bfseries, anchor=south west] at (0.05,5.05)
  {(b) Cluster-structured (ours)};
\draw[dashed, black!45] (3.45,0) .. controls (3.2,1.8) and (3.6,3.4) .. (3.3,5);
\draw[dashed, black!45] (3.4,2.15) .. controls (4.8,2.3) and (5.8,2.0) .. (7,2.3);
\node[cent] at (1.9,2.7) {};
\node[cent] at (5.2,3.7) {};
\node[cent] at (5.3,1.0) {};
\node[anchdot] at (2.0,2.4) {};
\node[posdot]  at (2.35,2.75) {};
\foreach \p in {(1.45,2.0),(2.65,1.85),(1.6,3.25),(2.85,2.95)}
  \node[negfill] at \p {};
\node[note, red!60!black] at (1.75,1.15) {50\% same cluster\\(hardest)};
\foreach \p in {(4.35,3.4),(5.0,4.25),(4.7,2.75)}
  \node[negfill] at \p {};
\node[note, red!60!black] at (4.75,4.75) {30\% nearest\\cluster};
\node[negfill] at (6.3,0.65) {};
\node[note, red!60!black] at (5.45,0.4) {20\% random};
\node[graydot] (cp) at (2.5,2.15) {};
\draw[red!70!black, thick] (2.36,2.01) -- (2.64,2.29);
\draw[red!70!black, thick] (2.36,2.29) -- (2.64,2.01);
\node[note, black!60, anchor=west] at (3.0,1.45)
  {co-positive: excluded};
\draw[-Stealth, black!50] (3.05,1.6) -- (2.62,2.05);
\end{scope}

\begin{scope}[shift={(0.6,-0.7)}]
\node[anchdot]  at (0,0) {};   \node[note, anchor=west] at (0.15,0) {anchor};
\node[posdot]   at (1.7,0) {}; \node[note, anchor=west] at (1.85,0) {positive};
\node[negopen]  at (3.6,0) {}; \node[note, anchor=west] at (3.75,0) {in-batch neg.};
\node[negfill]  at (6.1,0) {}; \node[note, anchor=west] at (6.25,0) {mined neg.};
\node[cent]     at (8.4,0) {}; \node[note, anchor=west] at (8.55,0) {centroid};
\end{scope}

\end{tikzpicture}%
}
\caption{In-batch InfoNCE draws negatives uniformly over the dataset, so
the anchor's actual look-alike neighborhood (dashed) is rarely sampled
(a). Our mining instead draws 50/30/20\% of negatives from the anchor's
own, nearest, and random clusters ranked by descriptor similarity, with
co-positive exclusion (b); because clusters are rebuilt from the current
descriptor, this pool self-corrects as training progresses.}
\label{fig:mining_comparison}
\end{figure}

\subsection{Training Objective}
Paired in-batch InfoNCE aligns anchor-positive descriptors
$\mathbf{q}_i,\mathbf{p}_i$,
\begin{equation}
\mathcal{L}_{align} = -\frac{1}{B}\sum_i \log
\frac{\exp(\mathbf{q}_i^T\mathbf{p}_i/\tau)}
{\sum_j \exp(\mathbf{q}_i^T\mathbf{p}_j/\tau)}, \quad \tau=0.10,
\end{equation}
and for sampled negatives $\mathbf{n}_{ij}$ from Section~\ref{sec:hard_mining}
a margin loss suppresses confusable candidates,
\begin{equation}
\mathcal{L}_{neg} = \mathbb{E}_{i,j}\Big[\max\big(0,
\tfrac{1}{2}(\mathbf{q}_i^T\mathbf{n}_{ij}+\mathbf{p}_i^T\mathbf{n}_{ij})
-0.60\big)^2\Big].
\end{equation}
A per-dimension variance term $\mathcal{L}_{var}$ prevents descriptor
collapse, and a soft cluster-consistency term $\mathcal{L}_{consistency}$
(symmetric KL between the anchor's and positive's softmax similarity to
all centroids) encourages paired crops to agree on cluster assignment.
The full objective is
\begin{equation}
\mathcal{L} = 1.0\,\mathcal{L}_{align} + 1.0\,\mathcal{L}_{var}
+ 0.2\,\mathcal{L}_{neg} + 2.0\,\mathcal{L}_{consistency}.
\end{equation}

\section{Experiments}
\subsection{Setup}
We evaluate on five UAV-satellite benchmarks, each with a model trained
or fine-tuned only on that benchmark's own official training split:
DenseUAV (2,256 locations, cross-altitude multi-positive pairing,
$\tau{=}30$\,m weighting) \cite{dai2024denseuav}; SUES-200 (official
120/80 split, 4 flight heights) \cite{zhu2023sues}; University-1652 (701
buildings) \cite{zheng2020university1652}; GTA-UAV same-area, loaded in
the official \texttt{pos\_semipos} mode (123,095 IoU-weighted pairs)
\cite{ji2024game4loc}; and UAV-VisLoc, processed with the Game4Loc tiling
protocol \cite{xu2024uavvisloc}. DenseUAV/SUES-200/University-1652 are
evaluated with each benchmark's own official closed-gallery evaluator;
GTA-UAV/UAV-VisLoc use the Game4Loc IoU-based open-map protocol
(positive IoU $\geq 0.39$). $K{=}8$--$16$ clusters, $K_c{=}2$ searched at
evaluation.

\subsection{Ablation: Does Cluster-Structured Mining Help?}
Table~\ref{tab:ablation} isolates the mining mechanism itself on
UAV-VisLoc (seq.~04) and GTA-UAV, holding the backbone fixed. Moving from
a flat descriptor to two-stage cluster-routed scoring, and then to
cluster-structured negative mining (cluster-Mahalanobis scoring),
recovers the errors concentrated in look-alike neighborhoods: R@1 rises
from 0.39 (flat, untrained) to 0.92 on UAV-VisLoc, and the same
cluster-routed descriptor matches or exceeds Game4Loc's own trained
model on GTA-UAV despite a simpler, single-descriptor design.

\begin{table}[t]
\caption{Ablation on UAV-VisLoc seq.~04 and GTA-UAV same-area. Bold =
best per block.}
\label{tab:ablation}
\centering
\footnotesize
\setlength{\tabcolsep}{3pt}
\begin{tabular}{>{\raggedright\arraybackslash}p{0.42\columnwidth}
                >{\centering\arraybackslash}p{0.12\columnwidth}
                >{\centering\arraybackslash}p{0.12\columnwidth}
                >{\centering\arraybackslash}p{0.15\columnwidth}}
\toprule
Variant & R@1 & R@5 & AP \\
\midrule
\multicolumn{4}{l}{\textit{UAV-VisLoc, seq.~04}} \\
\midrule
Flat, untrained backbone      & 0.392 & 0.602 & --    \\
Flat learned descriptor       & 0.623 & 0.817 & 0.301 \\
Two-stage, cluster-centered   & 0.637 & 0.968 & 0.529 \\
Two-stage, cluster-Mahalanobis (ours) & \textbf{0.919} & \textbf{1.000} & \textbf{0.719} \\
\midrule
\multicolumn{4}{l}{\textit{GTA-UAV same-area}} \\
\midrule
Game4Loc~\cite{ji2024game4loc} (paper) & 0.850 & 0.976 & \textbf{0.902} \\
Two-stage, cosine (ours) & \textbf{0.901} & 0.983 & -- \\
\bottomrule
\end{tabular}
\end{table}

\subsection{Comparison with State of the Art}
Tables~\ref{tab:soa_sues200}--\ref{tab:soa_denseuav} compare against the
strongest recent published methods on the three closed-gallery
benchmarks, all of which add part-partitioning heads (LPN, FSRA, MCCG,
CAMP), explicit spatial-domain alignment losses (DAC, MEAN), or
multi-directional attention modules (CEUSP) on top of the base
contrastive objective. Our single cluster-routed descriptor, adapted only
through cluster-structured mining, is competitive with or exceeds all of
them on at least one metric per benchmark, with no additional trainable
head or loss term. On DenseUAV, an optional post-hoc re-ranking of the
coarse top-10 using a quantized spatial-layout code (144 bytes/image, no
retraining) further separates near-duplicate overlapping tiles, raising
ConvNeXt-T from 86.14 to 94.38 R@1 -- exceeding CEUSP's published 89.45.

\begin{table}[t]
\caption{SUES-200 (Recall@1, \%, drone$\to$satellite, official 120/80
split, 80-tile test gallery). Bold = best per column.}
\label{tab:soa_sues200}
\centering
\footnotesize
\setlength{\tabcolsep}{4pt}
\begin{tabular}{l l cccc}
\toprule
Method & Backbone & 150\,m & 200\,m & 250\,m & 300\,m \\
\midrule
LPN~\cite{wang2021lpn}                 & ResNet-50  & 58.20 & 69.60 & 75.60 & 78.50 \\
FSRA~\cite{dai2022fsra}                & ViT-S      & 78.70 & 85.65 & 88.95 & 91.50 \\
MCCG~\cite{shen2024mccg}               & ConvNeXt-T & 79.96 & 85.01 & 90.47 & 94.20 \\
Sample4Geo~\cite{deuser2023sample4geo} & ConvNeXt-B & 94.75 & 96.75 & 97.25 & 97.20 \\
MEAN~\cite{chen2025mean}               & ConvNeXt-T & 95.50 & \textbf{98.38} & 98.95 & 99.17 \\
DAC~\cite{xia2024dac}                  & ConvNeXt-B & \textbf{96.80} & 97.48 & 98.20 & 97.58 \\
\midrule
\textbf{Ours} & Swin-B     & 96.30 & 98.25 & \textbf{99.00} & 99.15 \\
\textbf{Ours} & ConvNeXt-T & 91.77 & 95.20 & 96.90 & 97.52 \\
\textbf{Ours} & ConvNeXt-B & 94.97 & 97.28 & 98.50 & \textbf{99.35} \\
\bottomrule
\end{tabular}
\end{table}

\begin{table}[t]
\caption{University-1652 (R@1/AP, \%). D$\to$S: drone query, satellite
gallery. S$\to$D: satellite query, drone gallery. Bold = best per column.}
\label{tab:soa_u1652}
\centering
\footnotesize
\setlength{\tabcolsep}{3pt}
\begin{tabular}{l l cc cc}
\toprule
& & \multicolumn{2}{c}{D $\to$ S} & \multicolumn{2}{c}{S $\to$ D} \\
Method & Backbone & R@1 & AP & R@1 & AP \\
\midrule
LPN~\cite{wang2021lpn}                 & ResNet-50  & 75.93 & 79.14 & 86.45 & 74.79 \\
FSRA~\cite{dai2022fsra}                & ViT-S      & 84.51 & 86.71 & 88.45 & 83.37 \\
MCCG~\cite{shen2024mccg}               & ConvNeXt-T & 89.64 & 91.32 & 94.30 & 89.39 \\
Sample4Geo~\cite{deuser2023sample4geo} & ConvNeXt-B & 92.65 & 93.81 & 95.14 & 91.39 \\
MEAN~\cite{chen2025mean}               & ConvNeXt-T & 93.55 & 94.53 & 96.01 & 92.08 \\
DAC~\cite{xia2024dac}                  & ConvNeXt-B & \textbf{94.67} & \textbf{95.50} & 96.43 & 93.79 \\
\midrule
\textbf{Ours} & Swin-B     & 94.52 & 95.44 & \textbf{96.86} & \textbf{93.83} \\
\textbf{Ours} & ConvNeXt-T & 82.08 & 84.69 & 92.44 & 80.38 \\
\textbf{Ours} & ConvNeXt-B & 90.91 & 92.44 & 95.72 & 90.32 \\
\bottomrule
\end{tabular}
\end{table}

\begin{table}[t]
\caption{DenseUAV (drone$\to$satellite, same-area, full 18{,}198-tile
gallery). SDM@1 is DenseUAV's spatial-distance metric. Bold = best per
column. \textsuperscript{$\triangle$}MCCG row is our own reproduction from
official code (199 epochs, paper's recipe), not a secondhand figure -- see
extended version for details.}
\label{tab:soa_denseuav}
\centering
\footnotesize
\setlength{\tabcolsep}{3pt}
\begin{tabular}{l l ccc}
\toprule
Method & Backbone & R@1 & R@5 & SDM@1 \\
\midrule
DenseUAV baseline~\cite{dai2024denseuav} & ViT-S      & 83.01 & 95.58 & 86.50 \\
Sample4Geo~\cite{deuser2023sample4geo}   & ConvNeXt-B & 49.38 & 78.29 & 61.72 \\
MCCG~\cite{shen2024mccg}\textsuperscript{$\triangle$} & ConvNeXt-T & 90.26 & \textbf{97.90} & \textbf{91.82} \\
CAMP~\cite{wu2024camp}                   & ConvNeXt-B & 88.72 & --    & --    \\
CEUSP~\cite{xu2025ceusp}                 & ConvNeXt-T & 89.45 & 96.05 & 91.01 \\
DINOv2-GLFA~\cite{yang2025dinov2}         & DINOv2-B   & 86.27 & 96.83 & 88.87 \\
\midrule
\textbf{Ours}              & Swin-B     & 86.62 & 96.83 & 89.50 \\
\textbf{Ours}              & ConvNeXt-T & 86.14 & 96.91 & 89.12 \\
\textbf{Ours}              & ConvNeXt-B & 87.22 & 96.95 & 89.85 \\
\textbf{Ours} + VQ re-rank & ConvNeXt-T & \textbf{94.38} & \textbf{98.24} & -- \\
\bottomrule
\end{tabular}
\end{table}

\section{Conclusion}
Restructuring the negative pool around the same K-means partition used
for routing -- rather than adding a part-supervision head or a
spatial-alignment loss -- is enough to make a single pooled descriptor
competitive with, and often better than, recent multi-component UAV
cross-view methods. Because the partition is rebuilt from the current
descriptor every epoch, the mining pool self-corrects toward whichever
confusions the model has not yet resolved, without any additional
trainable parameters. Full per-altitude/per-gallery breakdowns, the
vector-quantized re-ranking derivation, and results on UAV-VisLoc/GTA-UAV
appear in the extended version of this work.

\bibliographystyle{IEEEtran}
\bibliography{references}

\end{document}
